Behavior-Driven Development (BDD) plays a vital role in software testing by bridging the gap between business requirements and automated tests. However, manual creation of Gherkin scenarios and Behave step definitions is labor-intensive and prone to human error. While Large Language Models (LLMs) show significant potential in code generation, there is a lack of rigorous empirical evaluation comparing different prompting techniques for joint scenario and step definition generation using open-source models. In this study, we evaluate Zero-Shot, Few-Shot, and Chain-of-Thought (CoT) prompting techniques using the Qwen2.5-7B-Instruct model on a dataset of $N=100$ user stories. Our results show that both Few-Shot (mean Cosine Similarity = $0.8141$, $p = 0.0007$) and CoT ($0.8100$, $p = 0.0046$) prompting significantly exceed the $0.80$ quality threshold, with Few-Shot performing significantly better than Zero-Shot ($p = 0.0034$) while maintaining a perfect $100\%$ syntax execution rate. These findings suggest that Few-Shot prompting provides the most balanced and resource-efficient approach for practical BDD automation pipelines.
